
\textbf{Proposition 3.11.} \uline{Soient} @@M0@@ \uline{un} @@M1@@-anneau, \uline{et} @@M2@@, @@M3@@ \uline{des éléments de} @@M4@@ \uline{tels que} @@M5@@ \uline{pour} @@M6@@. \uline{On pose} @@M7@@~; \uline{alors} :
\begin{enumerate}[label=\roman*)]
\item \uline{pour} @@M8@@, on a @@M9@@
\item \uline{pour} @@M10@@, @@M11@@ \uline{est la} @@M12@@-ième \uline{fonction symétrique des} @@M13@@.
\end{enumerate}
Soit @@M14@@. On a alors @@M15@@. Comme @@M16@@ pour @@M17@@, et @@M18@@ pour @@M19@@, @@M20@@. La proposition résulte alors de ce que @@M21@@.

\section*{4. @@M22@@-anneaux engendrés par générateurs et relations}

On se place toujours sur @@M23@@.

4.1. Soit @@M24@@ un @@M25@@-anneau. On considère une famille d'indéterminées @@M26@@, et soit @@M27@@~; on veut définir sur @@M28@@ une structure de @@M29@@-anneau prolongeant celle de @@M30@@.

Pour se donner un homomorphisme @@M31@@, il suffit de se donner les images des @@M32@@, @@M33@@, ce qui définit un homomorphisme d'anneaux : @@M34@@, et pour qu'il définisse une structure de @@M35@@-anneau sur @@M36@@, il faut vérifier la commutativité du diagramme (3.7.1). On pose donc :
@@M37@@
Le diagramme (3.7.1) étant un diagramme d'homomorphismes d'anneaux, il suffit de vérifier sa commutativité pour les @@M38@@, @@M39@@ (elle est connue par hypothèse sur les éléments de @@M40@@), ce qui résulte facilement des définitions (4.1.1) et de ce que @@M41@@ et @@M42@@ sont des @@M43@@-homomorphismes.

